CUDA Fortranでは以下のように記述するとカーネルが呼び出される．

\begin{lstlisting}[language=Fortran]
call calc_keep_x<<<blocks, threads, stream>>>(nx, ny, nz, ruvwp, T, E)
\end{lstlisting}

三重の不等号はシェブロンと呼ばれる．カーネルを呼ぶにはブロック数，スレッド数，実行するストリームを指定する．デフォルトのストリームは0番である．ストリーム番号を変えるとカーネルが並行実行されることがある（計算リソースが不足するため，実際に並行実行されることはほとんどない）．

blocksとthreadsは以下のように宣言する．

\begin{lstlisting}[language=Fortran]
type(dim3), parameter :: threadsE = dim3(32, 2, )
type(dim3), parameter :: blocksE  = dim3((nx-1+threadsE%x-1)/threadsE%x, &
                                         (ny-1+threadsE%y-1)/threadsE%y, &
                                         (nz-1+threadsE%z-1)/threadsE%z)
\end{lstlisting}

\noindent
シェブロンで呼び出されるsubroutineにはglobal属性を付与する．

\begin{lstlisting}[language=Fortran]
attributes(global) subrouinte calc_keep_x(nx, ny, nz, ruvwp, T, E)
\end{lstlisting}

\noindent
global属性が付与されたsubroutine内ではblockIdx, blockDim, threadIdxが使用できる．

\begin{itemize}
\item blockIdx:現在実行されているblockの番号
\item blockDim:blockに含まれるthreadの数
\item threadIdx:block内でのthread番号（最大値はblockDim）
\end{itemize}

\noindent
これらをi, j, kと対応させるには以下のように書く．

\begin{lstlisting}[language=Fortran]
i = (blockIdx%x-1)*blockDim%x + threadIdx%x
j = (blockIdx%y-1)*blockDim%y + threadIdx%y
k = (blockIdx%z-1)*blockDim%z + threadIdx%z
\end{lstlisting}
